% unidades/10-familia.tex
\chapter{👨‍👩‍👧‍👦 家族 — Familia}
\unidadheader{10}{家族}{Familia}{Hablar de los miembros de la familia y contar personas}

\begin{tcolorbox}[vocabbox, title={📌 Vocabulario Nuevo}]
\begin{tabularx}{\linewidth}{llX}
\toprule
\textbf{Japonés} & \textbf{Español} & \textbf{Tipo} \\
\midrule
\vocabitem{父}{ちち}{mi padre}{sust.}
\vocabitem{母}{はは}{mi madre}{sust.}
\vocabitem{兄}{あに}{mi hermano mayor}{sust.}
\vocabitem{姉}{あね}{mi hermana mayor}{sust.}
\vocabitem{弟}{おとうと}{mi hermano menor}{sust.}
\vocabitem{妹}{いもうと}{mi hermana menor}{sust.}
\vocabitem{お父さん}{おとうさん}{padre (de otro)}{sust.}
\vocabitem{お母さん}{おかあさん}{madre (de otro)}{sust.}
\vocabitem{〜人}{にん}{contador de personas}{contador}
\vocabitem{一人}{ひとり}{una persona}{sust.}
\vocabitem{二人}{ふたり}{dos personas}{sust.}
\bottomrule
\end{tabularx}
\end{tcolorbox}

\subsection*{🗣️ Frases Clave}

\frase{\ruby{家族}{かぞく}は\ruby{何人}{なんにん}ですか？}{¿Cuántas personas hay en tu familia?}
\frase{\ruby{五人}{ごにん}\ruby{家族}{かぞく}です。}{Somos una familia de cinco.}
\frase{\ruby{兄}{あに}が\ruby{二人}{ふたり}います。}{Tengo dos hermanos mayores.}
\frase{お\ruby{姉}{ね}さんはどこにいますか？}{¿Dónde está tu hermana mayor?}

\subsection*{⚙️ Gramática}

\patron{Contar personas}
  {数字 + 人 (にん)}
  {Excepciones: 一人 (ひとり), 二人 (ふたり), 四人 (よにん)}

\patron{Existencia de personas}
  {〜に 人 が います}
  {\ej{\ruby{教室}{きょうしつ}に\ruby{学生}{がくせい}が\ruby{十人}{じゅうにん}います。}{Hay 10 estudiantes en el salón.}
  \ej{\ruby{家}{いえ}に\ruby{父}{ちち}がいます。}{Mi padre está en casa.}}

\begin{tcolorbox}[ejerciciobox, title={📝 Ejercicios Rápidos}]
\begin{enumerate}
  \item ¿Cómo dices "Tengo una hermana menor"?
  \item ¿Cómo dices "Hay 4 personas"?
  \item Traduce: \ruby{田中}{たなか}さんのご\ruby{家族}{かぞく}はどこにいますか。
\end{enumerate}
\vspace{0.4em}
\textbf{Respuestas:}
1) \ruby{妹}{いもうと}がいます。\quad
2) \ruby{四人}{よにん}います。\quad
3) ¿Dónde está la familia del Sr. Tanaka?
\end{tcolorbox}

\begin{tcolorbox}[kanjibox, title={🀄 Kanjis de la Unidad}]
\begin{tabularx}{\linewidth}{cllXX}
\toprule
\textbf{Kanji} & \textbf{On} & \textbf{Kun} & \textbf{Significado} & \textbf{Vocab} \\
\midrule
\kanjirow{父}{ふ}{ちち}{padre}{\ruby{父}{ちち} / \ruby{祖父}{そふ}}
\kanjirow{母}{ぼ}{はは}{madre}{\ruby{母}{はは} / \ruby{祖母}{そぼ}}
\kanjirow{兄}{けい・きょう}{あに}{hermano mayor}{\ruby{兄}{あに} / \ruby{兄弟}{きょうだい}}
\kanjirow{姉}{し}{あね}{hermana mayor}{\ruby{姉}{あね} / \ruby{姉妹}{しまい}}
\kanjirow{弟}{てい・だい}{おとうと}{hermano menor}{\ruby{弟}{おとうと} / \ruby{兄弟}{きょうだい}}
\kanjirow{妹}{まい}{いもうと}{hermana menor}{\ruby{妹}{いもうと} / \ruby{姉妹}{しまい}}
\bottomrule
\end{tabularx}
\end{tcolorbox}
